\chapter{Artin transport, norms and specialization}
\label{v4-arith:at:chapter}

The arithmetic fusion category has two kinds of transport.
An Artin automorphism moves its marked junctions from one lattice
to another. A finite lattice covering instead multiplies
arithmetic values over periodic orbits. We construct both maps
and compare their tensor properties.

The distinction already appears for \(\varepsilon=2+\sqrt3\)
and \(I=\mathbb Z[\sqrt3]\). The category at this period has
simples \(1,g,\rho\), with
\[
 g^2=1,\qquad g\rho=\rho g=\rho,\qquad
 \rho^2=1+g+2\rho.
\]

Changing to period \(\varepsilon^2\) gives a positive map of
fusion rings that sends the new noninvertible simple to
\(1+g+4\rho\). The unit occurs once in this image.
Tensorator coherence would force a commutation relation on that
one-dimensional space which two nonzero scalars cannot satisfy.
The ring map survives. Its linear tensor lift does not.

We first transport the full marked category under reciprocity.
Periodic points then organize multiplicative norms and their
parameter algebras. The unit multiplicity gives obstructions to
tensor lifts of every even period ratio. A further cocycle
calculation excludes period three for the lattice above.
Finally, norms of lines and indexed tensors provide positive
multiplicative constructions, with their own coefficient domains.

\section{Reciprocity on the marked family}
\label{v4-arith:at:family}

Use the arithmetic data of Chapter~\ref{v4-arith:mf:chapter}.
Thus \(K\) is a real quadratic field, \(I\subset K\) is an
oriented rank-two lattice, and \(\varepsilon>1\) is a totally
positive norm-one unit preserving \(I\). Write
\[
 G_I=(\varepsilon-1)^{-1}I/I,\qquad
 N=|G_I|=\varepsilon+\varepsilon^{-1}-2.
\]
Let \(m\) be the exponent of \(G_I\), and write
\(\mathcal O=(I:I)=\mathbb Z+c_0\mathcal O_K\) for its multiplier
order, with conductor \(c_0\).

The normalized arithmetic values are
\(e_{I,g}=E_{I,\varepsilon}(g)/E_{I,\varepsilon}(0)\).
They lie in the narrow ray field \(L=H^+_{\mathcal O,m}\).
We use arithmetic Artin symbols: at an unramified prime their
reduction acts by the residue-field norm power.

The reciprocity input is Theorem~2 of
\emph{Stark units for real quadratic fields and reciprocity laws}.
It is stated in \eqref{v4-arith:mf:artin-input}.
The coefficients \(b_{I,g}=b_0e_{I,g}\), with
\(b_0=(\varepsilon-1)/N\), define the Leavitt algebra \(B_I\)
and the marked category \(\mathcal C_I\).
The algebra isomorphism below transports all four generator
types \(X_g,X'_g,T_g,T'_g\), including their algebraic duals.



Now let \(\mathfrak a\) be an integral invertible
\(\mathcal O\)-ideal prime to \(c_0m\), and put \(J=\mathfrak a^{-1}I\).
The natural label map is
\begin{equation}\label{v4-arith:at:eq:label-map}
 \phi_{\mathfrak a}:G_I\longrightarrow G_J,\qquad
 x+I\longmapsto x+J.
\end{equation}

\begin{lemma}\label{v4-arith:at:lem:labels}
The map \(\phi_{\mathfrak a}\) is an isomorphism, and \(J\) has
multiplier order \(\mathcal O\).
\end{lemma}
\begin{proof}
Coprimality gives \(z\in\mathfrak a\) with
\(z\equiv1\pmod{m\mathcal O}\). The integer
\(h=\operatorname{Nm}(z)=z\bar z\) belongs to \(\mathfrak a\)
and is prime to \(m\). Therefore \(h(J/I)=0\).
The kernel of \(\phi_{\mathfrak a}\) is killed by both \(h\)
and \(m\), so is zero. Both groups have order \(N\), proving
surjectivity. If \(tJ\subseteq J\), multiplication by the invertible
ideal \(\mathfrak a\) gives \(tI\subseteq I\), and the converse
follows by multiplication by \(\mathfrak a^{-1}\).
\end{proof}

Extend coefficients to \(L\), and write
\(\sigma_{\mathfrak a}=\operatorname{Art}(\mathfrak a)|_L\).
Since \(b_0\in K\), the arithmetic law gives
\[
 \sigma_{\mathfrak a}(b_{I,g})
       =b_{J,\phi_{\mathfrak a}g}.
\]
All reconstruction coefficients are rational expressions in the
\(b_g\) and \(N\). Thus \(Q,\chi,a,c_1,\beta\) transform by the
same label map, while \(d\) is fixed.

\begin{theorem}\label{v4-arith:at:thm:artin}
There is a \(\sigma_{\mathfrak a}\)-semilinear strong tensor
equivalence
\[
 \Phi_{\mathfrak a,I}:\mathcal C_{I,L}
       \longrightarrow\mathcal C_{\mathfrak a^{-1}I,L}
\]
preserving \(\rho\), sending \(\alpha_g\) to
\(\alpha_{\phi_{\mathfrak a}g}\), and preserving the specified
duality maps and normalized networks. For two such ideals,
\[
 \Phi_{\mathfrak b,\mathfrak a^{-1}I}\Phi_{\mathfrak a,I}
       =\Phi_{\mathfrak a\mathfrak b,I}
\]
on the explicit models.
\end{theorem}
\begin{proof}
Apply \(\sigma_{\mathfrak a}\) to scalars and \(\phi_{\mathfrak a}\)
to every generator label. This defines a semilinear algebra
isomorphism \(\Theta_{\mathfrak a}:B_{I,L}\to
B_{J,L}\). The coefficient formulas show
\(\Theta_{\mathfrak a}\rho_I=\rho_J\Theta_{\mathfrak a}\) and the
corresponding equality for \(\alpha_g\).
Thus intertwiners map to intertwiners, and
\eqref{v4-arith:mf:tensor-map} is preserved. Finite direct sums
and idempotent splittings are preserved as well.

The zero label and \(d\) are fixed, proving duality compatibility.
The normalization of Proposition~\ref{v4-arith:mf:network-value} is preserved.
Finally, Artin symbols multiply and both composed label maps send
a representative to its class in \((\mathfrak a\mathfrak b)^{-1}I\).
The two semilinear maps therefore agree on scalars and generators.
\end{proof}

This is transport on the lattice-indexed family. Identifying its
targets with one fixed category would require lattice identifications.
This construction lifts the arithmetic law
\eqref{v4-arith:mf:artin-input} to the marked category.


The scalar normalization matters for the unnormalized values.
Choose an extension of \(\sigma_{\mathfrak a}\) to
\(L(\sqrt\varepsilon)\). For this extension, reciprocity reads
\[
 \sigma_{\mathfrak a}\bigl(E_{I,\varepsilon}(x)\bigr)
 =\frac{\sigma_{\mathfrak a}(\sqrt\varepsilon)}
        {\sqrt\varepsilon}\,
   E_{\mathfrak a^{-1}I,\varepsilon}(x).
\]
The possible sign disappears for \(e_{I,x}\).
It cannot be dropped from a statement about the raw values.

\subsection{Transport of singular interfaces and finite symmetry}

Let \(A\) be the level-one Heisenberg cut cobar over \(K\)
used in Chapter~\ref{v4-arith:asi:chapter}. Put \(R_I=B_I^{\mathrm{op}}\).
For an endomorphism \(\eta\) of \(B_I\), let \(P_\eta\) have
underlying space \(B_I\) and actions
\[
 a^{\mathrm{op}}b\,c^{\mathrm{op}}=\eta(c)ba.
\]
The relative realization is
\[
 \mathfrak R_I=A_L\otimes_LR_{I,L},\qquad
 \mathfrak P_\eta=A_L\otimes_LP_{\eta,L}.
\]
The map \(\Theta_{\mathfrak a}\) transports these coefficients.
The system \(A\) is defined over \(K\), so its coordinates,
translation, residues, local vacuum and cut maps are fixed.

Let \(P,Q\) be composable coefficient bimodules over
\(R=R_{I,L}\). In bar degree \(n\), the comparison
\[
 (A_L\otimes P)\star^{\mathbb L}_{A_L\otimes R}
                  (A_L\otimes Q)
 \simeq A_L\otimes(P\otimes_R^{\mathbb L}Q)
\]
multiplies the \(n+2\) geometric factors in order.
It leaves the coefficient bar
\(P\otimes R^{\otimes n}\otimes Q\).
Thus it commutes with \(\Theta_{\mathfrak a}\) in every degree.
The common multibar proves the same assertion for every
associativity and unit constraint.

Naturality also transports the connected comparison
\eqref{v4-arith:asi:theta} and its weight-one interface
\eqref{v4-arith:asi:tangent-comparison}.
All normal derivatives and spectator states are retained.

The finite receiver has an equally explicit transport.
Take split finite semisimple rigid categories over
characteristic-zero fields \(k,k'\), with simple units.
Let \(F:\mathcal C\to\mathcal C'\) be a specified
\(\sigma:k\to k'\)-semilinear tensor equivalence.
Choose its identifications of simple objects, with permutation
\(\pi\) of their labels. They give semilinear maps
\[
 U_i:P_{S_i}\longrightarrow P'_{S'_{\pi(i)}}
\]
compatible with the tensorators, units and duality pairings of
Section~\ref{v4-arith:asi:weak-hopf}.
The reconstructed algebras are
\(H_{\mathcal C}=\bigoplus_i\operatorname{End}_k(P_{S_i})\)
and \(H_{\mathcal C'}\).

\begin{proposition}\label{v4-arith:at:weak-transport}
The formula
\[
 \Theta_\sigma(h)_{\pi(i)}=U_i h_iU_i^{-1}
\]
defines a semilinear weak Hopf isomorphism.
These maps compose with the marked equivalences.
\end{proposition}
\begin{proof}
Conjugation preserves multiplication.
The compatibility of \(U\) with the balanced tensorator and
matching-endpoint projection preserves the coproduct
\eqref{v4-arith:asi:weak-coproduct}.
Compatibility with the unit and its coordinate-sum functional
preserves the counit. Compatibility with the duality pairings
preserves their transposes, hence the antipode.
Composing the maps \(U\) gives the stated composition law.
\end{proof}

In matrices after scalar twisting, this formula is
\(U_i\,\sigma(h_i)\,U_i^{-1}\), where \(U_i\) now denotes the
linear matrix of the twisted map.
The Artin equivalences supply the required identifications:
they preserve \(X_0,dX'_0\) and every marked network.
The assertion therefore includes the weak unit, categorical
junctions and arithmetic markings.


\section{Coherence on a fixed presentation}
\label{v4-arith:at:fixed-presentation}

Transport of a marked algebra has two distinct forms. Base change
moves the whole presentation to its conjugate. Transport on one fixed
presentation first requires identifications with those conjugates,
and their associativity has an obstruction. We compute this
obstruction before discussing multiplicative norms, whose coherence uses different data.

Let $k$ be a field and let $B$ be an associative unital $k$-algebra
with specified generators, relations and marked intertwiners. Include
all roots and normalization constants in $k$. An automorphism
$\sigma$ of $k$ carries these data to their coefficientwise
conjugates. Polynomial relations and compositions are preserved,
so this construction gives coherent semilinear transport on the
family of conjugate presentations. Identifying its value with a
prescribed arithmetic relabelling requires maps of the marked
presentations themselves.

\subsection{The central associativity class}

Regard a group $\Gamma$ acting on $k$ as a discrete group. Suppose
we are given $\sigma$-semilinear algebra automorphisms $\phi_\sigma$
of $B$ and units $u_{\sigma,\tau}\in B^\times$ satisfying
\begin{equation}\label{v4-arith:at:eq:weak}
 \phi_\sigma\phi_\tau
  =\operatorname{Ad}(u_{\sigma,\tau})\phi_{\sigma\tau},
 \qquad \phi_1=1,\qquad u_{1,\sigma}=u_{\sigma,1}=1.
\end{equation}
The last two unit equalities are part of the normalization.
Set
\begin{equation}\label{v4-arith:at:eq:cocycle}
 c_{\sigma,\tau,\nu}
 =(u_{\sigma,\tau}u_{\sigma\tau,\nu})^{-1}
    \phi_\sigma(u_{\tau,\nu})u_{\sigma,\tau\nu}.
\end{equation}
The restrictions of the $\phi_\sigma$ to $Z(B)$ define an honest
$\Gamma$ action, because inner automorphisms are trivial on the center.
Write \(\sigma(z)=\phi_\sigma(z)\) for \(z\in Z(B)\).

\begin{proposition}[Associativity of a fixed presentation]
\label{v4-arith:at:prop:associator}
The elements \eqref{v4-arith:at:eq:cocycle} form a normalized $3$-cocycle
with values in this $\Gamma$-module $Z(B)^\times$. Replacing
$u_{\sigma,\tau}$ by $z_{\sigma,\tau}u_{\sigma,\tau}$ for a
normalized central $2$-cochain changes $c$ by
\begin{equation}\label{v4-arith:at:eq:coboundary}
 (\delta z)_{\sigma,\tau,\nu}
  =\frac{\sigma(z_{\tau,\nu})z_{\sigma,\tau\nu}}
          {z_{\sigma,\tau}z_{\sigma\tau,\nu}}.
\end{equation}
For the supplied automorphisms and inner comparisons, coherent
implementers with $c=1$ exist exactly when
$[c]=0$ in $H^3(\Gamma,Z(B)^\times)$.
\end{proposition}
\begin{proof}
Both $u_{\sigma,\tau}u_{\sigma\tau,\nu}$ and
$\phi_\sigma(u_{\tau,\nu})u_{\sigma,\tau\nu}$ implement the
same comparison between
$\phi_\sigma\phi_\tau\phi_\nu$ and $\phi_{\sigma\tau\nu}$.
Their ratio is therefore central.

Compare the five
parenthesizations of four automorphisms using these implementers.
The central triple ratios commute with every remaining factor.
The two paths around the pentagon give
\[
 \sigma(c_{\tau,\nu,\kappa})
 c_{\sigma,\tau\nu,\kappa}c_{\sigma,\tau,\nu}
 =c_{\sigma\tau,\nu,\kappa}c_{\sigma,\tau,\nu\kappa}.
\]
This is the multiplicative cocycle equation.

The unit equalities
in \eqref{v4-arith:at:eq:weak} make $c$ equal to $1$ when any argument is
the identity. Substitution of $zu$ into
\eqref{v4-arith:at:eq:cocycle} gives \eqref{v4-arith:at:eq:coboundary}.
Finally, two units implementing the same inner automorphism differ
by a central unit. These cochains exhaust the choices for fixed
\(\phi\). Their coboundaries therefore give exactly the stated
vanishing criterion.
\end{proof}

When $c=1$, the direct sum
\begin{equation}\label{v4-arith:at:eq:crossed}
\begin{gathered}
 B\rtimes_{\phi,u}\Gamma=\bigoplus_{\sigma\in\Gamma}BU_\sigma,
 \\
 (bU_\sigma)(dU_\tau)
   =b\phi_\sigma(d)u_{\sigma,\tau}U_{\sigma\tau}
\end{gathered}
\end{equation}
is a unital associative algebra over $k^\Gamma$.
Indeed expansion of a triple product uses
$\phi_\sigma\phi_\tau=\operatorname{Ad}(u_{\sigma,\tau})
\phi_{\sigma\tau}$ to reduce its two parenthesizations to the
two products of implementers in \eqref{v4-arith:at:eq:cocycle}.
For general $c$, the right-associated product is the left-associated
one multiplied in its coefficient by $c_{\sigma,\tau,\nu}$.
This also makes the defect of associativity explicit.

Normalization of $\phi_1$ alone is insufficient. For example, take
$B=\mathbb Q$, $\Gamma=\mathbb Z/3$ written additively, all
$\phi=1$, and choose $u_{0,1}=2$, $u_{0,2}=3$, with all other
$u=1$. The inner relation still holds, but
$c_{0,1,1}=3/2$. The units at an identity argument are central and
can be normalized first. The displayed \(c\) becomes normalized
only after this correction.

The strict maps on an arithmetic family in
Theorem~\ref{v4-arith:at:thm:artin} and the fixed-presentation datum
\eqref{v4-arith:at:eq:weak} have different domains. A family action does
not by itself supply the identifications and inner implementers
needed to define this class. Conversely, the class does not
obstruct coefficientwise transport on the whole conjugate family.

\section{Periodic points and normalized norms}
\label{v4-arith:at:sec:norms}

The lattice operations have a geometric organization independent of
the chiral background. Form the real two-torus
\[
 \mathbb T_I=(K\otimes_{\mathbb Q}\mathbb R)/I.
\]
Multiplication by \(\varepsilon\) induces a hyperbolic automorphism
\(u_I\) with eigenvalues \(\varepsilon,\varepsilon^{-1}\).
For every integer \(r\geq1\),
\begin{equation}\label{v4-arith:at:eq:periodic}
 \operatorname{Fix}(u_I^r)
       =(\varepsilon^r-1)^{-1}I/I,\qquad
 N_r=|\operatorname{Fix}(u_I^r)|
       =\operatorname{Tr}(\varepsilon^r)-2.
\end{equation}
Indeed, a lift is fixed precisely when
\((\varepsilon^r-1)x\in I\). The real-linear map
\(\varepsilon^r-1\) is invertible, and its determinant gives
the stated cardinality.

The formal periodic-point zeta function is
\[
 \exp\left(\sum_{r\geq1}\frac{N_r}{r}t^r\right)
       =\frac{(1-t)^2}{1-\operatorname{Tr}(\varepsilon)t+t^2}.
\]
This follows by expanding the logarithms of
\((1-\varepsilon t)\), \((1-\varepsilon^{-1}t)\) and \((1-t)^2\).

An inclusion \(I\subset J\) of lattices preserved by
\(\varepsilon\) gives an equivariant finite covering
\(f:\mathbb T_I\to\mathbb T_J\). A fiber over a point fixed by
\(u_J^r\) is permuted by \(u_I^r\).
Its elements need not all have period \(r\).
The orbit lengths determine the norm.

More generally, let \(u\) denote compatible automorphisms on spaces
with finite equivariant coverings. A multiplicative periodic datum
with values in a field \(M\) consists of
\(w(x,r)\in M^\times\), for \(r\geq1\) and \(u^rx=x\), satisfying
\begin{equation}\label{v4-arith:at:eq:period-datum}
 w(ux,r)=w(x,r),\qquad w(x,rs)=w(x,r)^s\quad(s\geq1).
\end{equation}
For an equivariant finite covering \(f\), define
\begin{equation}\label{v4-arith:at:eq:orbit-norm}
 (f_*w)(z,r)=
 \prod_{\mathcal O\in u^r\backslash f^{-1}(z)}
       w(y_{\mathcal O},r|\mathcal O|),
\end{equation}
where \(y_{\mathcal O}\) is any representative of the orbit.
The first identity in \eqref{v4-arith:at:eq:period-datum} makes the choice
irrelevant.

\begin{proposition}\label{v4-arith:at:prop:norms}
The datum \(f_*w\) satisfies \eqref{v4-arith:at:eq:period-datum}.
For composable equivariant finite coverings,
\((g\circ f)_*w=g_*(f_*w)\).
If the covers have constant degree and preserve a distinguished
fixed point \(0\), the normalized operations
\begin{equation}\label{v4-arith:at:eq:normalized-norm}
 N_f^0(w)(z,r)=\frac{(f_*w)(z,r)}{(f_*w)(0,r)}
\end{equation}
also satisfy \(N_g^0N_f^0=N_{g\circ f}^0\).
\end{proposition}
\begin{proof}
Applying \(u\) permutes the orbits in a fiber and preserves their
weights. To verify the power law, an orbit of length \(s\) under
\(u^r\) splits into \(\gcd(s,t)\) orbits under \(u^{rt}\).
Each new weight is the old weight raised to
\(t/\gcd(s,t)\). Their product is its \(t\)-th power.

For transitivity, fix a point of period dividing \(r\) in the
last space. An orbit in the middle fiber has length \(s\).
Above one of its representatives, an orbit under \(u^{rs}\)
has some length \(t\). It gives one orbit of length \(st\)
in the composite fiber, and every composite orbit is obtained
uniquely this way. Both products therefore have the same factors
\(w(x,rst)\).

Put \(v=f_*w\). Before applying \(g_*\), normalization divides
the factor on an orbit of length \(s\) by
\(v(0,rs)=v(0,r)^s\). The total denominator is
\(v(0,r)^{\deg(g)}\), independent of the target point.
It cancels when the result is normalized at zero.
Ordinary transitivity then proves normalized transitivity.
\end{proof}

For the arithmetic system take \(M=\overline{\mathbb Q}\) and
\[
 w_I(x,r)=\frac{E_{I,\varepsilon^r}(x)}
                    {E_{I,\varepsilon^r}(0)}.
\]
The period identities give \eqref{v4-arith:at:eq:period-datum}.
Proposition 3 of \emph{Stark units for real quadratic fields and
reciprocity laws} gives the product over the orbits in a covering
fiber. Dividing that identity by its value at zero removes its
root-of-unity multiplier and proves
\begin{equation}\label{v4-arith:at:eq:arithmetic-norm}
 N_f^0(w_I)=w_J.
\end{equation}
The product of the orbits' zero-value normalization factors is
\(\varepsilon^{r\deg(f)/2}\), independent of their decomposition.
Thus replacing raw values by \(w_I\) gives the same normalized law.

If \(H\subset G_{I,\varepsilon}\) and \(J/I=H\), the law reduces,
for \(x\in G_{I,\varepsilon}\), to
\begin{equation}\label{v4-arith:at:eq:subgroup-norm}
 e_{J,\varepsilon}(\pi x)
       =\prod_{h\in H}\frac{e_{I,\varepsilon}(x+h)}
                              {e_{I,\varepsilon}(h)}.
\end{equation}
The denominator is the zero-point normalization. The image of
\(G_I/H\) need not fill \(G_J\). The other points are handled by
the longer orbits in \eqref{v4-arith:at:eq:orbit-norm}.

One can encode each weight as a one-dimensional \(M\)-space with
its scalar automorphism. The products then become tensor products
of lines, and normalization tensors with the dual reference line.
This realizes the norm law on lines with automorphisms. It does
not construct tensor functors between the fusion categories.



For example, take \(J=(\varepsilon-1)^{-1}I\).
Then \(G_I=J/I\), so its image in \(G_J\) consists only of zero.
But \(|G_J|=N\). If \(N>1\), every other fixed point of
\(\mathbb T_J\) has a fiber with no fixed lift.
The longer orbits in \eqref{v4-arith:at:eq:orbit-norm} are necessary.

The arithmetic identity uses quotient coverings of lattices.
It is not invariant under every equivariant map of tori.
For the order-two example of \eqref{v4-arith:asi:example-unit},
translation by the nonzero fixed point exchanges the values
\(1,e_1\). Its normalized norm sends \(e_1\) to
\(e_1^{-1}\ne e_1\).

\subsection{The algebra of normalized periodic parameters}


These nonlinear operations become algebra maps on parameter spaces.
Define the Laurent algebra
\[
 \mathcal P_I=K[z_{x,r}^{\pm1}:r\ge1, x\in\operatorname{Fix}(u^r)]
 \big/(z_{ux,r}-z_{x,r}, z_{x,rs}-z_{x,r}^s, z_{0,r}-1),
\]
where the power relation is imposed only when $u^rx=x$.
The pullback for $f$ is the finite monomial formula
\begin{equation}\label{v4-arith:at:norm-algebra}
 f^*(z_{z,r})=
 \frac{\prod_{O\subset f^{-1}(z)}z_{y_O,r|O|}}
      {\prod_{O\subset f^{-1}(0)}z_{y_O,r|O|}}.
\end{equation}

The orbit argument proves that it respects the relations and that these
pullbacks compose. Applying Theorem~\ref{v4-arith:asi:all-arity} over $K$ gives
all-arity chiral algebra maps
$\mathcal A_{\mathcal P_J}\to\mathcal A_{\mathcal P_I}$.
Arithmetic evaluation to $\overline K$ commutes with them.
For any finite collection of periods, Artin transport is compatible
under the common coefficient extension and its stated coprimality
conditions. A single finite ray class field is not asserted for all
periods at once.


\section{Changing the period}
\label{v4-arith:at:sec:period}

Write \(G_r=(\varepsilon^r-1)^{-1}I/I\).
The quotient map \(I/(\varepsilon^r-1)I\to I/(\varepsilon-1)I\)
gives a canonical surjection
\begin{equation}\label{v4-arith:at:eq:period-label}
 \pi_r:G_r\longrightarrow G_1,\qquad
 x\longmapsto(1+\varepsilon+\cdots+\varepsilon^{r-1})x.
\end{equation}
This is distinct from the inclusion \(G_1\subset G_r\) used in
the arithmetic power identity.

Let \(R_{G_1}\) be the integral fusion ring with basis
\(\alpha_g,\rho\) and rules \eqref{v4-arith:mf:fusion-rule}.
Put
\[
 A_G=\sum_{g\in G_1}\alpha_g,\qquad W_r=(1+\rho)^r-1.
\]
The element \(W_r\) is positive in the fusion basis, since it is
a sum of positive tensor powers with nonnegative integer coefficients.

\begin{theorem}\label{v4-arith:at:thm:period-ring}
There is a positive unital ring homomorphism
\[
 \Psi_r:R_{G_r}\longrightarrow R_{G_1},\qquad
 \alpha_h\longmapsto\alpha_{\pi_r(h)},\qquad
 \rho_r\longmapsto W_r.
\]
These homomorphisms compose under successive changes of period.
\end{theorem}
\begin{proof}
In \(R_{G_1}\otimes\mathbb Q\), the element \(e=A_G/N\)
satisfies \(e^2=e\), \(e\rho=\rho\), and \((1-e)\rho=0\).
In the component with unit \(e\), set \(z=e+\rho\).
The fusion rule gives
\[
 z^2-(N+2)z+e=0,\qquad z+z^{-1}=(N+2)e.
\]
The second equality uses the inverse \((N+2)e-z\).
The recurrence for \(z^r+z^{-r}\), together with
\(N+2=\operatorname{Tr}(\varepsilon)\), yields
\(z^r+z^{-r}=\operatorname{Tr}(\varepsilon^r)e\).

Since \(1+\rho=(1-e)+z\) has orthogonal components,
\(W_r=z^r-e\). Consequently
\begin{equation}\label{v4-arith:at:eq:period-square}
 W_r^2=N_rW_r+(N_r/N)A_G.
\end{equation}
The integer \(N_r/N\) is the size of \(\ker\pi_r\). Thus
\[
 \sum_{h\in G_r}\alpha_{\pi_r(h)}=(N_r/N)A_G,
\]
and \eqref{v4-arith:at:eq:period-square} is exactly the image of the
source fusion relation. The relations involving invertibles follow
from the group map and \(\alpha_gW_r=W_r\).

The rational calculation proves an integral identity because the
fusion basis is free over \(\mathbb Z\).

For successive periods \(r,s\), the label multipliers multiply
to \(1+\varepsilon+\cdots+\varepsilon^{rs-1}\). On the other
generator the composite is
\((1+W_r)^s-1=(1+\rho)^{rs}-1\).
This proves composition on every basis element.
\end{proof}

Both roots of the dimension equation are respected:
\[
 W_r(\varepsilon-1)=\varepsilon^r-1,\qquad
 W_r(\varepsilon^{-1}-1)=\varepsilon^{-r}-1.
\]
For \(K=\mathbb Q(\sqrt5)\), \(I=\mathcal O_K\), and
\(\varepsilon=(3+\sqrt5)/2\), one has \(N=1\),
\(N_2=5\), and \(G_2\simeq\mathbb Z/5\).
The map sends every invertible simple at period two to the unit
and sends \(\rho_2\) to \(1+3\rho\). Indeed,
\[
 (1+3\rho)^2=10+15\rho=5+5(1+3\rho).
\]



A tensor lift requires maps on multiplicity spaces.
Choose the fusion bases of \eqref{v4-arith:mf:F} and let
\(\mathscr F\) be the additive functor specified by \(\Psi_r\).
Its tensorator would consist of invertible maps
\[
 J_{X,Y}:\mathscr F(X)\otimes\mathscr F(Y)
              \longrightarrow\mathscr F(X\otimes Y)
\]
satisfying
\begin{equation}
\begin{aligned}
 &\mathscr F(a_{X,Y,Z})J_{X\otimes Y,Z}(J_{X,Y}\otimes1)\\
 &\qquad=J_{X,Y\otimes Z}(1\otimes J_{Y,Z})
       a_{\mathscr F(X),\mathscr F(Y),\mathscr F(Z)}.
\end{aligned}
 \label{v4-arith:at:tensorator-equation}
\end{equation}
There are also the two unit equations.
The finite fusion bases make these polynomial equations with
invertibility conditions. The next obstruction uses the
tensorator only on invertible objects and the unit channel.


\section{Quotient maps and their multiplicities}\label{v4-arith:at:sec:norm}

Let $\mathcal C_N$ and $\mathcal C_M$ have the near-group rules above, with groups $G,H$ of
orders $N,M$, and let $\pi:G\twoheadrightarrow H$ have kernel of order $\ell$, so $N=\ell M$.
Suppose
an additive strong tensor functor $F$ sends $\alpha_g$ to $\alpha_{\pi(g)}$. Write $\tau$ for
the noninvertible simple in the target and $S=\sum_{h\in H}[\alpha_h]$. Since
$\alpha_g\rho=\rho$, the class of $F\rho$ is invariant under all translations of $H$. Thus
\[
 [F\rho]=aS+r[\tau],\qquad a,r\in\mathbb Z_{\ge0}.
\]

The two coefficients contain all possible multiplicities. No additional simple objects exist in
the specified target.

\begin{theorem}[Exact fusion-ring obstruction]\label{v4-arith:at:thm:norm-obstruction}

A positive unital ring homomorphism with the indicated map on invertible simples exists if and
only if there are integers $a\ge0$, $r>0$ satisfying
\begin{equation}\label{v4-arith:at:eq:norm-obstruction}
 r=\ell-2a,\qquad (M+4)a(\ell-a)=\ell(\ell-1).
\end{equation}
Consequently the existence of an additive strong tensor functor with that invertible-sector map
requires \eqref{v4-arith:at:eq:norm-obstruction}. The equations do not alone construct its tensorator.
\end{theorem}
\begin{proof}
In the target $S^2=MS$, $S\tau=M\tau$, and $\tau^2=S+M\tau$. Hence
\[
 (aS+r\tau)^2=(Ma^2+r^2)S+(2arM+r^2M)\tau.
\]
The image of the source relation is
\[
 \ell S+N(aS+r\tau).
\]

First $r\ne0$. Otherwise the positive number $aM$ would solve $d^2=Nd+N$, whose positive root
$(N+\sqrt{N^2+4N})/2$ is irrational: $N^2+4N$ lies strictly between $(N+1)^2$ and $(N+2)^2$ for
$N\ge1$. Comparing the $\tau$ coefficients and dividing by $rM$ gives $r=\ell-2a$. Substitution
in
the $S$ coefficient gives the second equation.

Conversely these computations verify every source
relation, including the group actions, so the assignment defines the stated ring homomorphism.
It is nonnegative and sends the unit to the unit.
\end{proof}

For a proper quotient, $a=0$ is impossible. For $\ell=2,3,4$, the equations and $r>0$ have no
solutions: at $\ell=2$ there is no positive $a<\ell/2$. At $\ell=3$ the only candidate is $a=1$,
which
would force $M=-1$. At $\ell=4$ it forces $M=0$. In particular, there is no such tensor functor
from the $N=4$ near-group category to the $M=2$ category along a quotient of their invertible
sectors. This obstructs the additive tensor functor itself.

There are nontrivial passing ring tests. For $a=1$, the equations give
\[
 M=\ell-4,\qquad N=\ell(\ell-4),\qquad r=\ell-2\qquad(\ell\ge5).
\]
Thus $N=12$, $M=2$, $\ell=6$ admits the positive ring assignment $\rho\mapsto S+4\tau$. This
passing ring test does not assert existence of every numerical near-group category, and no
tensor functor is inferred.

For two given categories, the remaining problem is a finite system
of tensorator equations on their multiplicity spaces. Those equations use both associators. The
finite receiver of Section~\ref{v4-arith:mf:finite} turns these into matrices, but
it does not construct a solution.

For the fixed target $M=2$, the criterion has a useful arithmetic form:
\[
 3r^2=\ell(\ell+2),\qquad (\ell+1)^2-3r^2=1,\qquad a=(\ell-r)/2.
\]
In particular, define integers $x_j,r_j$ by $x_j+r_j\sqrt3=(2+\sqrt3)^j$ and put $\ell_j=x_j-1$.
Then
\[
 (N,M,\ell,a,r)=(2\ell_j,2,\ell_j,(\ell_j-r_j)/2,r_j)
\]
is a passing based-ring datum for every $j\ge1$. The first proper cases are $(12,2,6,1,4)$ and
$(50,2,25,5,15)$.

The recurrence $x_{j+1}=2x_j+3r_j$, $r_{j+1}=x_j+2r_j$ preserves
$x_j^2-3r_j^2=1$. Reduction modulo two gives $\ell_j\equiv r_j$, and $\ell_j\ge r_j$ for
$j\ge1$.
These facts prove integrality and nonnegativity of $a$. These maps specialize the positive
period maps of Theorem~\ref{v4-arith:at:thm:period-ring} to the quadratic unit $2+\sqrt3$. They retain
the Pell constraint on the multiplicities.


An additional quick obstruction follows from Frobenius--Perron dimensions. A tensor functor
preserves this positive dimension character. Alternatively, this follows here directly by
applying the target positive dimension to the source relations. Write $d_N=(N+\sqrt{N^2+4N})/2$
and define $d_M$ similarly. They satisfy
\[
 d_N=aM+r d_M.
\]

Unequal quadratic fields $\mathbb Q(\sqrt{N^2+4N})$ and $\mathbb Q(\sqrt{M^2+4M})$ prohibit such
a map. Equality of these fields is necessary, not sufficient. Here the dimension condition
follows directly from the two fusion relations, without a separate categorical reconstruction.

The unit channel also constrains a lift when its multiplicity exceeds
one. We use the arithmetic endomorphism presentation of
\eqref{v4-arith:asi:fusion-identities}. Its source group is denoted by $G$,
its noninvertible object by $\rho$, and its nondegenerate bicharacter
by $\chi$. The intertwiners satisfy $U_gU_h=U_{g+h}$ by
\eqref{v4-arith:mf:U}. All functors below are linear over a common field
$F$ of characteristic zero. For a subgroup $H_0\subset G$, write
\[
 H_0^\perp=\{g\in G:\chi(g,h)=1\text{ for every }h\in H_0\}.
\]

\begin{proposition}[The full unit multiplicity]
\label{v4-arith:at:prop:unit-multiplicity}
Let $\mathsf F$ be a strong tensor functor and let $H_0\subset G$
be a subgroup whose invertible objects have image isomorphic to the
unit. Suppose
\[
 a=\dim_F\operatorname{Hom}(1,\mathsf F\rho)
\]
is finite and positive. Then
\begin{equation}\label{v4-arith:at:eq:unit-determinant}
 \chi(g,h)^a=1\quad(g,h\in H_0),\qquad
 aH_0\subset H_0^\perp.
\end{equation}
If $H_0$ has at most $d$ generators, then
\begin{equation}\label{v4-arith:at:eq:unit-divisibility}
 |H_0|\mid [G:H_0]a^d.
\end{equation}
\end{proposition}
\begin{proof}
Strictify the target, choose the unit identifications and put
\(M=\mathsf F\rho\). Write
\(J_{X,Y}:\mathsf F X\otimes\mathsf F Y\to\mathsf F(XY)\)
for the tensorator. For \(g,h\in H_0\), define invertible maps
\[
 A_g=J_{\alpha_g,\rho},\qquad
 B_h=J_{\rho,\alpha_h}^{-1}\mathsf F(U_h).
\]
Coherence on \((\alpha_g,\rho,\alpha_h)\) gives
\[
 J_{\alpha_g,\rho\alpha_h}
 =J_{\rho,\alpha_h}A_gJ_{\rho,\alpha_h}^{-1}.
\]
Naturality for \(U_h:\rho\to\rho\alpha_h\), together with
\(\alpha_g(U_h)=\chi(g,h)U_h\), then gives
\[
 A_gB_h=\chi(g,h)B_hA_g.
\]

These maps act by postcomposition on the entire space
\(\operatorname{Hom}(1,M)\).
Taking determinants gives \eqref{v4-arith:at:eq:unit-determinant}.

Write $H_0[a]$ for the kernel of
multiplication by $a$ on $H_0$. Its order divides $a^d$, by the cyclic
decomposition of $H_0$. Moreover $aH_0$ is a subgroup of $H_0^\perp$,
whose order is $[G:H_0]$ by nondegeneracy. The identity
$|H_0|=|H_0[a]|\,|aH_0|$ proves
\eqref{v4-arith:at:eq:unit-divisibility}.
\end{proof}

In particular, \(a=1\) forces \(H_0\subset H_0^\perp\).
For period doubling,
\[
 N_2=N(N+4),\qquad |\ker\pi_2|=N+4,\qquad
 W_2=S+(N+2)\rho.
\]
The unit multiplicity is one, so a tensor lift would require
\((N+4)^2\le N(N+4)\). Their difference is \(4(N+4)>0\).
This proves the obstruction in the opening example.
For the Fibonacci target, two nonzero scalars would have to
satisfy a nontrivial fifth-root commutation relation.

For period maps the kernel has at most two generators. The resulting
divisibility condition excludes every even ratio of periods.

\begin{theorem}[Even ratios of periods]
\label{v4-arith:at:thm:even-periods}
Let $I$ be a full rank-two lattice in a real quadratic field and let
$\varepsilon>1$ be a norm-one unit with $\varepsilon I=I$.
Put
\[
 G_r=(\varepsilon^r-1)^{-1}I/I,\qquad
 N_r=\varepsilon^r+\varepsilon^{-r}-2,
 \qquad N=N_1.
\]
Let $\pi_r:G_r\to G_1$ be multiplication by
$(\varepsilon^r-1)/(\varepsilon-1)$.
Assume that the arithmetic near-group categories for the periods in
question are given with the nondegenerate bicharacters and
endomorphism presentations above. If $r\ge2$ is even, the positive
period ring map of Theorem~\ref{v4-arith:at:thm:period-ring},
\[
 \alpha_g\longmapsto\alpha_{\pi_r(g)},\qquad
 \rho_r\longmapsto W_r=(1+\rho)^r-1,
\]
has no linear strong tensor lift. The same conclusion holds for
$G_{rs}\to G_s$ whenever the ratio $r$ is even, and after any
characteristic-zero extension of the coefficient field.
\end{theorem}
\begin{proof}
Put $t=N+2$, $V_j=\varepsilon^j+\varepsilon^{-j}$, and define
$u_{-1}=0$, $u_0=1$ and $u_j=tu_{j-1}-u_{j-2}$ for $j\ge1$.
In the target fusion ring write $S=\sum_{g\in G_1}\alpha_g$.
The relations $S^2=NS$, $S\rho=N\rho$ and $\rho^2=S+N\rho$ give
\begin{equation}\label{v4-arith:at:eq:period-multiplicity}
 W_r=a_rS+u_{r-1}\rho,
 \qquad Na_r=u_{r-1}-u_{r-2}-1.
\end{equation}
For example, multiplication of $W_r$ by $1+\rho$, followed by
addition of $\rho$, gives the formula for $W_{r+1}$. Thus all these
multiplicities are integers. Evaluating in the two roots
$\varepsilon,\varepsilon^{-1}$ also gives
\[
 Na_r=\frac{(\varepsilon^r-1)(\varepsilon^{r-1}-1)}
                  {\varepsilon^{r-1}(\varepsilon+1)}>0
 \qquad(r\ge2).
\]

The kernel $H_0$ of $\pi_r$ has order $N_r/N$ and index $N$ in
$G_r$. It is a subgroup of $(\mathbb Q/\mathbb Z)^2$, so it has at
most two generators. A lift would have unit multiplicity $a_r$.

Proposition~\ref{v4-arith:at:prop:unit-multiplicity} would therefore make
\begin{equation}\label{v4-arith:at:eq:period-quotient}
 \frac{N^2a_r^2}{N_r}
     =\frac{V_{r-1}-2}{N+4}
\end{equation}
an integer. The equality follows by inserting the preceding formula
and $N_r=(\varepsilon^r-1)^2/\varepsilon^r$.

But $V_0=2$, $V_1=t$ and $V_j=tV_{j-1}-V_{j-2}$. Modulo $N+4$,
this gives $V_j\equiv2(-1)^j$. For even $r$ the numerator on the
right of \eqref{v4-arith:at:eq:period-quotient} is congruent to $-4$.
Since $N+4\ge5$, it is not divisible by $N+4$.

Replacing $\varepsilon$ by $\varepsilon^s$ gives the assertion at
base period $s$. The unit multiplicity and the root-of-unity
commutator relations are unchanged by extension of coefficients.
\end{proof}

For an odd ratio $r=2k+1$, the same quotient is
\[
 \frac{V_{2k}-2}{N+4}=N u_{k-1}^2.
\]
Thus every odd ratio passes this determinant test. In the Pell family
above, $j=2m$ gives
\[
 \ell_j=6r_m^2,\qquad a_j=r_m(3r_m-x_m),\qquad
 \frac{2a_j^2}{\ell_j}=4r_m^2-2x_mr_m+\frac13,
\]
whereas $j=2m+1$ gives $\ell_j=(x_m+3r_m)^2$,
$a_j=r_m(x_m+3r_m)$ and $2a_j^2/\ell_j=2r_m^2$.
For the standard Pell lattice, the first odd ratio has a further
obstruction: coherence relates the two projective kernel actions.


\begin{lemma}[Reciprocal kernel cocycles]
\label{v4-arith:at:lem:kernel-cocycles}
In Proposition~\ref{v4-arith:at:prop:unit-multiplicity}, suppose the target
unit has endomorphism field $F$. Write
$J_{X,Y}:\mathsf F X\otimes\mathsf F Y\to\mathsf F(XY)$ for
the tensorator. Choose $\mathsf F\alpha_g=1$ for
$g\in H_0$ and write $c(g,h)\in F^\times$ for the tensorator on
$(\alpha_g,\alpha_h)$. Put
\[
 T_h=J_{\rho,\alpha_h},\qquad
 A_g=J_{\alpha_g,\rho},\qquad
 B_h=T_h^{-1}\mathsf F(U_h).
\]
Then $c$ is a scalar $2$-cocycle and
\begin{align}
 A_gA_h&=c(g,h)A_{g+h},\notag\\
 B_gB_h&=c(g,h)^{-1}B_{g+h},\notag\\
 A_gB_h&=\chi(g,h)B_hA_g.
 \label{v4-arith:at:eq:kernel-cocycles}
\end{align}
These identities hold on $\operatorname{Hom}(1,\mathsf F\rho)$.
\end{lemma}
\begin{proof}
Coherence on three invertible objects gives the cocycle identity.
Coherence on $(\alpha_g,\alpha_h,\rho)$ gives the first operator
identity, since $\alpha_h\rho=\rho$. For the second, the source has
\[
 (U_g\otimes1_{\alpha_h})U_h=U_{g+h}.
\]

Naturality and coherence give, respectively,
\[
 \mathsf F(U_g\otimes1)T_h
       =J_{\rho\alpha_g,\alpha_h}\mathsf F(U_g),\qquad
 J_{\rho\alpha_g,\alpha_h}T_g=c(g,h)T_{g+h}.
\]
Consequently $\mathsf F(U_{g+h})=c(g,h)T_{g+h}B_gB_h$.
Cancelling $T_{g+h}$ proves the second identity. The third is the
relation used in Proposition~\ref{v4-arith:at:prop:unit-multiplicity}.
Postcomposition restricts all three to the unit channel.
\end{proof}

\begin{proposition}[The third period of the standard Pell lattice]
\label{v4-arith:at:prop:third-period}
For $I=\mathbb Z[\sqrt3]$ and $\varepsilon=2+\sqrt3$, the positive
period map $G_3\to G_1$ has no linear strong tensor lift between the
specified arithmetic categories. The conclusion holds after any
characteristic-zero extension of coefficients.
\end{proposition}
\begin{proof}
Extend the field to an algebraic closure. Here $|G_3|=50$ and
$1+\varepsilon+\varepsilon^2=5\varepsilon$. Thus the kernel is
\[
 H_0=(\tfrac15 I)/I\simeq(\mathbb Z/5)^2,
\]
with basis $\varepsilon/5,1/5$. Formula
\eqref{v4-arith:at:eq:period-multiplicity} gives $W_3=5S+15\rho$, so a
lift would have a five-dimensional unit channel.

On this kernel the arithmetic Gaussian is
given in Radchenko--Wheeler, Section~4.4, p.~20:
\[
 Q\bigl((k\varepsilon+l)/5\bigr)
       =\zeta^{k^2-kl+l^2},\qquad \zeta=e^{2\pi i/5}.
\]
Hence $\chi(g,h)=\zeta^{g^{\mathsf t}Kh}$ with
$K=\left(\begin{smallmatrix}2&-1\\-1&2\end{smallmatrix}\right)$
over $\mathbb F_5$.

The commutator of a scalar cocycle on $H_0$ is
an alternating bicharacter. It therefore has the form
$c(g,h)c(h,g)^{-1}=\zeta^{v\det(g,h)}$ for some $v\in\mathbb F_5$.
By Lemma~\ref{v4-arith:at:lem:kernel-cocycles}, the operators $A_gB_h$
give a projective representation of $H_0\oplus H_0$ with alternating
commutator matrix
\[
 \Omega_v=\begin{pmatrix}vJ&K\\-K&-vJ\end{pmatrix},
 \qquad J=\begin{pmatrix}0&1\\-1&0\end{pmatrix}.
\]

Its Pfaffian is $-v^2-3$, which is nonzero in $\mathbb F_5$:
$2=-3$ is not a square. The commutator form is nondegenerate.

Choose a symplectic basis for this four-dimensional form. The
corresponding operators form two independent Weyl pairs of order
five. Rescale their scalar fifth powers to one. The two commuting
clock operators then diagonalize, and the two shift operators carry
each joint eigenspace through twenty-five distinct weights of equal
dimension. Every nonzero representation consequently has dimension
divisible by twenty-five. This contradicts the five-dimensional
unit channel.
\end{proof}

\begin{question}\label{v4-arith:at:qn:tensorator}

For the arithmetic categories attached to $\varepsilon=2+\sqrt3$, which odd periods $j\ge3$ admit a
linear strong tensor lift of the positive period map? Its noninvertible image is
$a_jS+r_j\tau$, where $a_j=(\ell_j-r_j)/2$ and the integers $\ell_j,r_j$ are defined above.
Construct invertible maps on the finite multiplicity bimodules satisfying the unit and
associator equations, or exhibit an obstruction. For $I=\mathbb Z[\sqrt3]$,
Proposition~\ref{v4-arith:at:prop:third-period} makes $j=5$ the first remaining case.
\end{question}



\section{Multiplicative norms on coefficients}

The preceding obstructions concern additive strong tensor
functors between the specified fusion categories.
Two multiplicative constructions remain available.
The first acts on invertible modules.
The second tensors an indexed family of algebras and bimodules.

\subsection{Norms of invertible modules}

Let \(R\) be a commutative ring and let \(E/R\) be finite \'etale
of constant rank \(d\). For an invertible
$E$-module $L$ define the invertible $R$-module
\begin{equation}\label{v4-arith:at:eq:norm}
 N_{E/R}(L)=\det_R(L)\otimes_R\det_R(E)^{-1}.
\end{equation}
The source and target are the Picard groupoids of invertible
modules and their isomorphisms. Their monoidal operation is
tensor product.

\begin{theorem}[Tensor and tower coherence of the norm]
\label{v4-arith:at:thm:norm}
The construction \eqref{v4-arith:at:eq:norm} is strong symmetric monoidal:
\[
 N(L\otimes_EM)\simeq N(L)\otimes_RN(M),\qquad N(E)\simeq R.
\]
For finite etale towers $R\to E\to F$, it has coherent comparisons
$N_{F/R}\simeq N_{E/R}N_{F/E}$. All these maps commute with
arbitrary base change on this finite etale domain.
\end{theorem}
\begin{proof}
After an etale splitting, $E=R^d$ and
$L=(L_1,\ldots,L_d)$. Formula \eqref{v4-arith:at:eq:norm} becomes the
tensor product of the lines $L_i$. A permutation of sheets has
the same determinant sign in $\det_R(L)$ and $\det_R(E)$, so
the signs cancel. Thus the local tensor comparison is invariant
under every sheet permutation.

For a tower, split all sheets at once. Either norm is the tensor
product over the same final sheets, grouped first by their fibers
over $E$ or directly over $R$. Associativity and symmetry of tensor
products identify these expressions and every longer grouping.
They are invariant under changes of splitting and hence descend
by faithful flat descent. The same sheetwise description proves
base change and all tensor and tower coherences.
\end{proof}

On automorphisms of the unit line, this is
$u\mapsto\det_R(m_u)$, the multiplicative norm of units.
The determinant polynomial on all elements is not additive.
For \(E=R\times R\), it is \((a,b)\mapsto ab\).
The two orthogonal idempotents have norm zero while their sum
has norm one.
An additive functor on all modules cannot induce this map on
endomorphisms of the unit. A transfer on a Hall or fusion category
therefore needs its own construction or a specified multiplicative
subcategory.

For comparison, a finite scalar norm has an exact normalized tower
identity. If $\pi:G\twoheadrightarrow H$ is a surjection of finite
groups and $u$ has values in a commutative group of units, set
\[
 N_\pi(u)(h)=\prod_{\pi(g)=h}u(g),\qquad
 \mathcal N_\pi(u)(h)=N_\pi(u)(h)/N_\pi(u)(1_H).
\]
For $G\xrightarrow\pi H\xrightarrow\rho K$ one has
\begin{equation}\label{v4-arith:at:eq:normalized-tower}
 \mathcal N_\rho(\mathcal N_\pi(u))=\mathcal N_{\rho\pi}(u).
\end{equation}

Every fiber of \(\rho\) has cardinality \(|\ker\rho|\).
The intermediate denominator therefore occurs to the same power
in every fiber product. Final normalization cancels it, proving
\eqref{v4-arith:at:eq:normalized-tower}.
The group identity here is the reference point for normalization.

\subsection{The indexed tensor norm on coefficients}

An indexed tensor norm has a different target and is nonadditive.
Fix a field \(k\) and a finite set \(S\). For
dg algebras $B_s,E_s,D_s$ over $k$ and composable dg bimodules
${}_{B_s}P_s{}_{E_s}$ and ${}_{E_s}Q_s{}_{D_s}$, define
\[
 N_S(B_s)=\bigotimes_{s\in S}B_s,\qquad
 N_S(P_s)=\bigotimes_{s\in S}P_s.
\]
Fix an order temporarily to write formulas. The symmetric dg tensor comparison makes different
orders canonically equivalent, with their Koszul signs. The left and right actions on the tensor
of bimodules are the ordered tensor actions.

\begin{theorem}[Composition of indexed norms]\label{v4-arith:at:thm:indexed-norm}

There is a coherent derived comparison
\[
 \left(\bigotimes_sP_s\right)
 \otimes^{\mathbb L}_{\bigotimes_sE_s}
 \left(\bigotimes_sQ_s\right)
 \simeq\bigotimes_s(P_s\otimes^{\mathbb L}_{E_s}Q_s).
\]
It preserves regular bimodules, arbitrary finite towers of indexing sets, and permutations of
the sheets. It does not assert additivity in the repeated input $P$ of $P\mapsto P^{\otimes S}$.
\end{theorem}
\begin{proof}
Use free bar resolutions or $K$-flat representatives. On homogeneous tensors the comparison is
\[
 (\otimes_sp_s)\otimes(\otimes_sq_s)
 \longmapsto(-1)^{\epsilon}\otimes_s(p_s\otimes q_s),\qquad
 \epsilon=\sum_{s<t}|q_s||p_t|.
\]

Balancing on sheet $s$ becomes precisely balancing over $E_s$ after this permutation. The tensor
differential commutes with a Koszul permutation. The same permutation applied to the common
multibar compares the derived tensors. Its realizations agree because finite tensors over a
field preserve the required colimits and because geometric realizations over multiple simplicial
indices commute.

For three interfaces, the two routes make the same final permutation of all
factors. Thus their signs agree. Regular bimodules give the unit comparison. Grouping a finite
set into fibers and then grouping those fibers leaves the final tensor and its permutation
unchanged, proving tower coherence.
\end{proof}

Over a characteristic-zero field, this construction descends along a finite \'etale algebra.
Choose a splitting cover with sheet set $S$. The symmetric tensor action of sheet permutations
gives descent data for the algebras and bimodules.

For invertible modules in degree zero, it recovers the Picard
norm, with normalization
$\det(f_*L)\otimes\det(f_*\mathcal O)^{-1}$ for the finite \'etale map $f$. On
arbitrary objects it is polynomial rather than additive. For example,
\[
 (P\oplus Q)^{\otimes S}
 \cong\bigoplus_{U\subseteq S}
 \left(\bigotimes_{s\in U}P\right)\otimes
 \left(\bigotimes_{s\notin U}Q\right),
\]
with factors kept in sheet order. The mixed summands exhibit
the failure of additivity.

Let \(\mathscr A\) be a fixed singular background over \(k\),
as in Chapter~\ref{v4-arith:asi:chapter}.
To obtain interfaces, apply the relative realization after taking this
coefficient norm:
\[
 N_S(B_s)\longmapsto\mathscr A\otimes N_S(B_s),\qquad
 N_S(P_s)\longmapsto\mathscr A\otimes N_S(P_s).
\]
Theorem~\ref{v4-arith:at:thm:indexed-norm} and Theorem~\ref{v4-thm:relative-composition} give its relative
interface composition. The construction retains one copy of
\(\mathscr A\). An identification of \(\mathscr A^{\otimes S}\)
with \(\mathscr A\) is not among its maps.

Sheet permutations act on the coefficient tensor.
They supply no exchange of positive geometric word factors.
The target is the relative envelope with the indexed coefficient
algebra. Its comparison with another period's arithmetic fusion
category is an additional construction.


\section{Integral transport and residue fields}
\label{v4-arith:at:integral}

The integral presentation of
Proposition~\ref{v4-arith:mf:integral} also transports.
For a fixed period, put \(\mathcal S=\mathcal O_L[1/N]\).
The field automorphism \(\sigma_{\mathfrak a}\) preserves this
ring, and all four generator formulas are defined over it.

\begin{proposition}\label{v4-arith:at:specialization}
The map \(\Theta_{\mathfrak a}\) restricts to a semilinear
isomorphism of the integral Leavitt presentations over
\(\mathcal S\). It preserves their endomorphisms, displayed
splittings, cups, caps and marked network equations.
For a prime ideal \(\mathfrak q\subset\mathcal S\), it induces
an isomorphism between the reductions at \(\mathfrak q\) and
\(\sigma_{\mathfrak a}(\mathfrak q)\), over the induced
residue-field isomorphism.
\end{proposition}
\begin{proof}
The map permutes the generator labels and conjugates their
coefficients. Each normalized coefficient is a unit in
\(\mathcal S\), and its conjugate is the corresponding
coefficient for \(\mathfrak a^{-1}I\).
The inverse label map and inverse field automorphism give the
inverse integral presentation map.

All the stated equations are identities between finite
expressions over \(\mathcal S\).
Their validity follows either from the defining relations or
from the free normal form of
Proposition~\ref{v4-arith:mf:integral}.
Reducing these expressions preserves every identity.
The semilinear map carries the ideal
\(\mathfrak q B_I\) onto
\(\sigma_{\mathfrak a}(\mathfrak q)B_{\mathfrak a^{-1}I}\),
so it descends to the claimed isomorphism.
\end{proof}

This statement concerns the presentations and their displayed
maps. Intertwiner spaces are kernels of linear equations.
Their dimensions can increase after reduction, so semisimplicity
and the full list of simple objects require a separate argument.
The localization excludes primes dividing \(N\).
An integral model at those primes requires an additional construction.

A finite collection of periods can be treated over a common
number field containing their coefficients and every chosen
normalization. Inverting the product of their \(N_r\) gives
a common ring for the displayed formulas.
The Artin ideals must satisfy the coprimality conditions for
each period. No single finite ray field is asserted for all
periods simultaneously.

\section{The remaining tensor comparison}

The constructions have distinct domains.
Reciprocity gives semilinear equivalences on a marked family.
Orbit norms give compatible maps of parameter algebras.
Picard norms and indexed tensor norms give multiplicative maps
on their specified coefficient categories.
The positive period ring map retains both dimension characters,
while its tensorator must solve
\eqref{v4-arith:at:tensorator-equation}.

For the standard lattice \(\mathbb Z[\sqrt3]\), every even
period ratio and period three are excluded by the preceding
proofs. At period five the proposed image is
\[
 W_5=76S+209\rho,\qquad N_5=722,\qquad
 |\ker\pi_5|=361.
\]
Its determinant test passes:
\(N^2\,76^2/722=32\).
The reciprocal kernel relations and the full associator equation
remain to be solved on these multiplicity spaces.
For other lattices, the restricted arithmetic Gaussian must
also be determined before applying the period-three argument.

\begin{question}\label{v4-arith:at:nonlinear-comparison}
Construct a multiplicative correspondence between the marked
arithmetic fusion systems whose scalar evaluation is the
normalized orbit norm.
Specify its coefficient category, supported field operations,
unit and comparison with the indexed norm.
Does it preserve the arithmetic junctions and their full
collision diagrams?
\end{question}

The parameter-algebra maps already realize the scalar norm in
every arity. The question asks for its comparison with the
marked fusion objects and their fields.

\section{References}

\begin{itemize}
\item D.~Radchenko and C.~Wheeler.
\emph{Real quadratic fields and finite quantum dilogarithms. I}.
\href{https://arxiv.org/abs/2609.21892v2}{arXiv:2609.21892v2},
especially Section~4.4 and the fusion construction in Section~5.

\item \emph{Stark units for real quadratic fields and reciprocity laws}.
4 October 2026,
\href{https://www.danrad.net/papers/stark-reciprocity.pdf}
{the reciprocity and distribution paper},
Theorem~2, Propositions~2--3 and Appendix~A.
\end{itemize}
